% Part: computability
% Chapter: computability-theory
% Section: def-functions-self-reference

\documentclass[../../../include/open-logic-section]{subfiles}

\begin{document}

\olfileid{cmp}{thy}{slf}
\olsection{स्वसंदर्भ वापरून फलनांची व्याख्या}

फलनांची व्याख्या त्यांच्याच साहाय्याने करता येणे साधारणपणे
उपयुक्त असते. उदाहरणार्थ, संगणनक्षम फलने $k$, $l$
आणि~$m$ दिली असता, स्थिरबिंदू उपप्रमेयानुसार प्रत्येक~$y$
साठी पुढील समीकरण पूर्ण करणारे आंशिक संगणनक्षम फलन~$f$
अस्तित्वात असते:
\[
f(y) \simeq
\begin{cases}
k(y) & \text{जर $l(y) = 0$ असेल तर} \\
f(m(y)) & \text{अन्यथा.}
\end{cases}
\]
अधिक नेमके सांगायचे, तर
\[
g(x,y) \simeq
\begin{cases}
k(y) & \text{जर $l(y) = 0$ असेल तर} \\
\cfind{x}(m(y)) & \text{अन्यथा}
\end{cases}
\]
असे घेऊन, नंतर $\cfind{e}(y) = g(e,y)$ होईल असा
निर्देशांक~$e$ स्थिरबिंदू उपप्रमेयाने मिळवल्यावर~$f$ मिळते.

ठोस उदाहरण म्हणून, ``महत्तम साधारण विभाजक'' फलन
$\fn{gcd}(u,v)$ ची व्याख्या अशी करता येते:
\[
\fn{gcd}(u,v) \simeq
\begin{cases}
v & \text{जर $u = 0$ असेल तर} \\
\fn{gcd}(\fn{mod}(v, u), u) & \text{अन्यथा}
\end{cases}
\]
येथे $\fn{mod}(v, u)$ म्हणजे $v$ ला~$u$ ने भागल्यावर
उरणारी बाकी. स्थिरबिंदू उपप्रमेय वापरून $\fn{gcd}$
आंशिक संगणनक्षम आहे हे मिळते. (प्रत्यक्षात, $y$ मध्ये
$\tuple{u, v}$ या जोडीचा संकेतांक घेऊन हे वरच्या
रूपात मांडता येते.) मग~$u$ वरील निगमनाने $\fn{gcd}$
हे प्रत्यक्षात सर्वत्र परिभाषित आहे हे दिसते.

अर्थात, नुकत्याच पाहिलेल्या उदाहरणांपेक्षा बरीच अधिक
गुंतागुंतीची स्वसंदर्भात्मक व्याख्या करता येते. बहुतेक
प्रोग्रामिंग भाषा फलनांची व्याख्या त्यांच्याच साहाय्याने
करण्याचे कोणते ना कोणते साधन देतात. पण एखादे फलन
संगणित करणाऱ्या अल्गोरिदमच्या \emph{निर्देशांका}च्या
साहाय्याने त्याची व्याख्या करता येण्यापेक्षा हे काहीसे
कमी विलक्षण आहे. स्थिरबिंदू प्रमेय पूर्ण सर्वसाधारणतेने
नेमके तेच करू देते.

\end{document}
